\newpage
\chapter{General Decision-Making}

\section{Status of this Handbook}

This Handbook is a guideline setting out how Deans carry out their work, based on Chapter 1, Article 5 and Chapter 5, Article 1, paragraph 2 of the Academy Regulations.
Give priority to the Academy Charter and the Academy Regulations; do not use this Handbook alone to impose new obligations on Students.
If a passage conflicts with the Regulations or an announced practice, check the higher-ranking document and discuss a correction with the other Deans.

\section{Principles of the Role}

Deans support an environment in which Students can continue their inquiry activities (Chapter 5 of the Academy Regulations).
Treat Students as equal members, and do not judge the relative merits of activities by their subject, the person's experience, or their achievements.
Intervene in Students' activities only where grounds specified in the Regulations exist, and only to the minimum extent necessary. When several means are available, choose those that are less restrictive and easier to reverse or recover from.
For decisions affecting rights or interests, make the grounds and established facts clear so that you can explain the reasons to the person concerned.
Record decisions and procedures so that a successor can take over the work.

\section{Decision Procedure and Consultation}

When asked to act, check the applicable Charter and Regulations provisions, announced guidance, and decisions in similar cases.
If facts are missing, ask for further information. Unless the situation is urgent, explain that you will respond after checking.
The Dean responsible may handle routine administrative work. Before reaching a conclusion, however, share with the other Deans any decision affecting Students' rights or interests, any decision without precedent, or any decision that will become a standard for the Academy's operation.

\section{Explanation, Records, and Handover}

The person making a decision records its subject, the established facts, the applicable provisions, the people consulted, the conclusion, and communications with the person concerned and their timing.
Do not finalize a decision affecting rights or interests solely through mechanical processing; the person responsible must be able to explain the reasons.
Keep only information needed for the work and do not share records with people who are not involved.
When responsibility changes, hand over unfinished matters and the reasons for decisions to the successor.
